% SPDX-License-Identifier: Apache-2.0
\section{Introduction}
\label{sec:e-intro}

\IEEEPARstart{A}{} companion paper~\cite{merge} took two hardware
description languages that had been developed independently for one
project, showed that they did not compete, and merged them. LHDL described
dataflow and left time to the compiler. TxHDL described transactions over
buses and let a module block for as many cycles as a protocol took. Four
questions needed a decision about meaning, and the analysis behind
them~\cite{unification} resolved all four.

The merged language was then given a surface syntax biased toward
Rust~\cite{syntax}, construct by construct, because Rust is familiar and
because following it removed six defects that the two sources had rather
than restating them.

This paper asks the question that arrangement invites. If the syntax is
already Rust, why is there a separate language at all?

\subsection{The experiment}

Take the merged specification and delete every construct that Rust already
provides. Rename every construct whose name collides with a Rust concept,
so that nothing means two things. Keep what is left, and see how small it
is.

The answer is that seventeen of the twenty non-Rust constructs go, and what
remains is a library rather than a keyword list.

Two of the deletions pay for the exercise on their own, and
Sections~\ref{sec:e-time} and~\ref{sec:e-pipeline} state them. A third
result is negative and matters as much: Section~\ref{sec:e-walls} gives the
one thing an embedded language cannot do, which is the thing a separate
parser was buying.

\subsection{What here is Chisel, and what is not}

Half of this paper rebuilds Chisel~\cite{bachrach2012} in Rust, and it
should say so before the reader works it out. A unit as a struct,
interfaces with directed roles, \code{mux} and \code{with!}, generators
through host-language parameters, and elaboration by running the program:
that is Chisel's architecture, and Sections~\ref{sec:e-units}
to~\ref{sec:e-binding} restate it. What Rust adds there is that one
driver per wire becomes a compile-time move error where Chisel raises an
elaboration-time exception. That is a better error, earlier, and not a
different design.

The other half is not Chisel. Chisel has no \code{async}, no
\code{await} and no sequencing; every multi-cycle protocol in it is a
state machine written by hand, which is the thing both source languages
existed to remove. Sections~\ref{sec:e-time} and~\ref{sec:e-pipeline},
where an \code{async fn} is a sequence, a pipeline and a process
according to how it is driven, and where an operator's latency is the
mapping's decision rather than the designer's, are the contribution.
Their nearest ancestors are XLS~\cite{xls}, whose DSLX is a
Rust-shaped language that the compiler schedules into pipelines, and
Calyx~\cite{calyx} and Filament~\cite{filament}, not Chisel;
Section~\ref{sec:e-related} places them, and says why XLS is also the
counterexample.

That half was, when this paper was first written, the half that did
not exist. Every probe type-checked; none \term{lowered} an
\code{async fn}, which is to say none translated it into the
registers, wires and control that a synthesis tool accepts. Lowering
is the compiler's word for going down the abstraction ladder, and it
is what turns a value that is alive across an \code{.await} into a
register that the source never named. Section~\ref{sec:e-next} states
the experiment that was to settle whether the contribution is real,
and what it found when it was run: an \code{async fn} lowers to a
pipeline, a unit's \code{run} loop to a module, and a RISC-V
core~\cite{vreteno} to a netlist that simulates against the trace of
its own source.

\subsection{Everything here was compiled}

The first draft of this work argued about what Rust would accept. That was
replaced. A Rust toolchain is now part of the repository's hermetic build,
and every claim in this paper is a probe that was compiled, one per
question. Section~\ref{sec:e-probes} reports what each one answered, and
keeps the ones that failed, because a failure is an answer.

Four of those compilations changed the design rather than confirming it,
and Section~\ref{sec:e-probes} gives all four. Two show the shape.
Arbitrary-width arithmetic does not work on stable Rust. A unit's parallel
processes cannot take a mutable borrow of the unit. Neither was visible
from the language reference; both were visible in ten lines of code.
